\chapter{Risk, Validation and Regulation}\label{ch:iv2:risk-validation-and-regulation}

\begin{ivtrains}
Value at risk and expected shortfall by hand and when they break, backtests and how much they can really tell you, stress tests, factor risk, the capital rules that a risk answer has to survive, and how to validate a model. Every portfolio here is invented. Quantiles of the normal: 2.326 at 99\% and 1.96 at 97.5\%.
\end{ivtrains}

\begin{question}[$\star$ \iqroles{Risk, QR, Trader} \iqfirm{any} \iqlevel{junior}]\label{iq:iv2:risk-validation-and-regulation:1}
You are long \$100 million of an equity index with a daily volatility of 1\%. What is the 99\% one-day value at risk, or VaR, under a normal model with zero mean? What does the square-root-of-time rule give for ten days?
\end{question}

\begin{question}[$\star$ \iqroles{Risk, QR, Analyst} \iqfirm{any} \iqlevel{junior}]\label{iq:iv2:risk-validation-and-regulation:2}
You hold \$50 million of one index with a daily volatility of 1\% and \$30 million of another with a daily volatility of 1.5\%. Their correlation is 0.8. What is the one-day 99\% VaR? How much does diversification save?
\end{question}

\begin{question}[$\star$ \iqroles{Risk, QR} \iqfirm{any} \iqlevel{junior}]\label{iq:iv2:risk-validation-and-regulation:3}
Define expected shortfall. For a normal distribution with mean 0 and standard deviation 1, what is it at 97.5\% and at 99\%? How does the 97.5\% figure compare with the 99\% VaR?
\end{question}

\begin{question}[$\star$ \iqroles{Risk, Analyst, QR} \iqfirm{any} \iqlevel{junior}]\label{iq:iv2:risk-validation-and-regulation:4}
Over 250 days, your five worst daily losses in millions were 8, 6, 5, 4 and 3.5. Estimate the 99\% historical VaR and the 99\% expected shortfall. Say which convention you use.
\end{question}

\begin{question}[$\star$ \iqroles{Risk, QR} \iqfirm{bank} \iqlevel{junior}]\label{iq:iv2:risk-validation-and-regulation:5}
A bank's 99\% one-day VaR had 7 exceptions in the past 250 days. Under the Basel traffic light, which zone is that, and what does it do to the multiplier of 3?
\end{question}

\begin{question}[$\star$ \iqroles{Risk, QR, QT} \iqfirm{any} \iqlevel{junior}]\label{iq:iv2:risk-validation-and-regulation:6}
An exponentially weighted volatility uses $\lambda = 0.94$. What is the half-life of a return's weight, in days? Yesterday's daily volatility estimate was 1\%, and today's return is 2\%. What is the new estimate?
\end{question}

\begin{question}[$\star$ \iqroles{Risk, Analyst, QR} \iqfirm{any} \iqlevel{junior}]\label{iq:iv2:risk-validation-and-regulation:7}
A stock has a beta of 1.2 to a market with 15\% volatility, and 20\% idiosyncratic volatility. What are its total volatility and its $R^2$? What is the volatility of an equal-weighted portfolio of 100 such stocks, with independent idiosyncratic parts?
\end{question}

\begin{question}[$\star$ \iqroles{Risk, Dev, QD} \iqfirm{any} \iqlevel{junior}]\label{iq:iv2:risk-validation-and-regulation:8}
An account's value goes 100, 120, 90, 130, 65, 80. What is the maximum drawdown? How do you compute it in one pass?
\end{question}

\begin{question}[$\star$ \iqroles{Risk, Trader, QT} \iqfirm{any} \iqlevel{junior}]\label{iq:iv2:risk-validation-and-regulation:9}
You are long a call with delta 0.5 and gamma 0.02 on a \$100 stock with a daily volatility of 2\%. What is the delta-normal one-day 99\% VaR per option? Is that too high or too low once you add gamma?
\end{question}

\begin{question}[$\star$ \iqroles{Risk, Trader, Analyst} \iqfirm{bank, hedge fund} \iqlevel{junior}]\label{iq:iv2:risk-validation-and-regulation:10}
A bond book has a DV01 of \$50,000, and the yield has a daily volatility of 6 basis points. What is the one-day 99\% VaR?
\end{question}

\begin{question}[$\star\star$ \iqroles{Risk, QR} \iqfirm{any} \iqlevel{junior, experienced}]\label{iq:iv2:risk-validation-and-regulation:11}
Two independent bonds each lose 100 with probability 4\% and nothing otherwise. What is the 95\% VaR of each one and of the pair? Compare the expected shortfall at 95\%. What property does this show that VaR lacks?
\end{question}

\begin{question}[$\star\star$ \iqroles{Risk, QR} \iqfirm{bank} \iqlevel{junior, experienced}]\label{iq:iv2:risk-validation-and-regulation:12}
Your 99\% VaR has 7 exceptions in 250 days. Run the Kupiec test. Do you reject at 5\%? What are the exact binomial tail probability and the test's verdict at 4 exceptions and at 0?
\end{question}

\begin{question}[$\star\star$ \iqroles{Risk, QR} \iqfirm{bank} \iqlevel{junior, experienced}]\label{iq:iv2:risk-validation-and-regulation:13}
Across 250 day-pairs, a 99\% VaR has these counts: no exception then no exception, 238; no exception then exception, 5; exception then no exception, 4; exception then exception, 3. Compute the transition probabilities and run Christoffersen's independence test.
\end{question}

\begin{question}[$\star\star$ \iqroles{Risk, QR} \iqfirm{any} \iqlevel{junior, experienced}]\label{iq:iv2:risk-validation-and-regulation:14}
Daily returns have a first-order autocorrelation of 0.2, following an AR(1). By how much does the square-root-of-time rule understate ten-day volatility? What happens with a negative autocorrelation?
\end{question}

\begin{question}[$\star\star$ \iqroles{Risk, QR} \iqfirm{any} \iqlevel{junior, experienced}]\label{iq:iv2:risk-validation-and-regulation:15}
Returns follow a Student-$t$ with 4 degrees of freedom, rescaled to a variance of 1. The 1\% quantile of the unscaled $t_4$ is $-3.747$. What are the 99\% VaR and expected shortfall, against the normal numbers?
\end{question}

\begin{question}[$\star\star$ \iqroles{Risk, QR, QD} \iqfirm{any} \iqlevel{junior, experienced}]\label{iq:iv2:risk-validation-and-regulation:16}
You estimate a 99\% VaR by Monte Carlo on a normal P\&L with 10,000 paths. What is the standard error of that quantile, in standard deviations and as a share of VaR? How many paths would bring it to 0.5\%?
\end{question}

\begin{question}[$\star\star$ \iqroles{Risk, QR} \iqfirm{any} \iqlevel{junior, experienced}]\label{iq:iv2:risk-validation-and-regulation:17}
For two positions with dollar volatilities of 0.5 and 0.45 and a correlation of 0.8, split the one-day 99\% VaR into contributions that add up to the total. Why do they add up? How do you confirm it by a bump?
\end{question}

\begin{question}[$\star\star$ \iqroles{Risk, Trader, Analyst} \iqfirm{bank, hedge fund} \iqlevel{junior, experienced}]\label{iq:iv2:risk-validation-and-regulation:18}
A book is long \$50 million of equities with beta 1, has a DV01 of \$40,000 long duration, and has a CS01 of \$30,000 (it loses that when spreads widen 1 basis point). Apply equities $-20\%$, rates $-100$ basis points and spreads $+150$. What is the P\&L? How much further must equities fall for a loss of \$15 million?
\end{question}

\begin{question}[$\star\star$ \iqroles{Risk, QR, Analyst} \iqfirm{hedge fund, bank} \iqlevel{junior, experienced}]\label{iq:iv2:risk-validation-and-regulation:19}
Three stocks have weights 0.5, 0.3 and 0.2, and betas on two factors of (1.2, 0.3), (0.8, $-0.2$) and (1.0, 0.5). The factor covariance has variances 0.04 and 0.01 and covariance 0.006. Idiosyncratic variances are 0.03, 0.02 and 0.04. What are the portfolio's volatility and the share of its variance that comes from factors?
\end{question}

\begin{question}[$\star\star$ \iqroles{Risk, Trader, QT} \iqfirm{hedge fund} \iqlevel{junior, experienced}]\label{iq:iv2:risk-validation-and-regulation:20}
You are long \$10 million of one stock and short \$10 million of a peer, each with a daily volatility of 2\%. What is the one-day 99\% VaR at a correlation of 0.9? What is it if the correlation falls to 0.5? What does that say about using recent correlation?
\end{question}

\begin{question}[$\star\star$ \iqroles{Risk, QR} \iqfirm{any} \iqlevel{junior, experienced}]\label{iq:iv2:risk-validation-and-regulation:21}
Distribution A is a standard normal. Distribution B is a standard normal with probability 99.5\%, and a loss of 50 with probability 0.5\%. Compare the 99\% VaR and the 99\% expected shortfall. Why have regulators moved towards shortfall? What do you lose by moving?
\end{question}

\begin{question}[$\star\star$ \iqroles{Risk, Trader, QT} \iqfirm{hedge fund, bank} \iqlevel{junior, experienced}]\label{iq:iv2:risk-validation-and-regulation:22}
You hold \$10 million of a stock with normal returns, a daily volatility of 2\% and a half-spread of 20 basis points, so a sale costs 20 basis points of the amount sold. You must sell in five equal parts, one at each day's close, so the whole position is at risk on the first day. What is the 99\% VaR of the sale, with the spread added as a sure cost? What is it if you can sell everything at the first close?
\end{question}

\begin{question}[$\star\star$ \iqroles{Risk, QR, QD} \iqfirm{bank} \iqlevel{junior, experienced}]\label{iq:iv2:risk-validation-and-regulation:23}
A desk asks you to validate its new pricing model before it goes live. What do you do? Give the steps in order, and what would make you refuse it.
\end{question}

\begin{question}[$\star\star$ \iqroles{Risk, QR} \iqfirm{bank} \iqlevel{junior, experienced}]\label{iq:iv2:risk-validation-and-regulation:24}
A 99\% VaR goes through a 250-day traffic light. If the model is right, what is the chance that it stays green? If the true exception rate is 2\% or 3\%, what is the chance that it stays green, and that it goes red?
\end{question}

\begin{question}[$\star\star$ \iqroles{Risk, Analyst, QR} \iqfirm{bank} \iqlevel{junior, experienced}]\label{iq:iv2:risk-validation-and-regulation:25}
Capital is $\max(\mathrm{VaR}_{t-1}, m\,\overline{\mathrm{VaR}}) + \max(\mathrm{sVaR}_{t-1}, m\,\overline{\mathrm{sVaR}})$. VaR was 12 yesterday and averages 10 over 60 days. Stressed VaR was 18 and averages 20, and $m = 3.65$. What is the capital? How high would yesterday's VaR need to be to matter? What does stressed VaR add?
\end{question}

\begin{question}[$\star\star\star$ \iqroles{Risk, QR, Trader} \iqfirm{any} \iqlevel{experienced}]\label{iq:iv2:risk-validation-and-regulation:26}
A delta-hedged book is short options with gamma $-500$ per dollar squared, and the stock's daily move is normal with a standard deviation of \$2. Ignore theta. What does delta-normal VaR say? What is the exact 99\% VaR? What does a mean-and-variance normal approximation give?
\end{question}

\begin{question}[$\star\star\star$ \iqroles{Risk, QR} \iqfirm{any} \iqlevel{experienced}]\label{iq:iv2:risk-validation-and-regulation:27}
Your returns have skewness $-0.5$ and excess kurtosis 3. Give the Cornish-Fisher 1\% quantile. Now use the formula on a Student-$t$ with 6 degrees of freedom, scaled to unit variance, whose 1\% quantile of the unscaled $t_6$ is $-3.143$. What do you learn?
\end{question}

\begin{question}[$\star\star\star$ \iqroles{Risk, QR} \iqfirm{bank} \iqlevel{experienced}]\label{iq:iv2:risk-validation-and-regulation:28}
A desk reports 97.5\% expected shortfall under a normal model. Over $T = 250$ days, a test statistic is $Z = 1 + \sum_t X_t I_t/(T \times 0.025 \times \mathrm{ES})$, with $X_t$ the P\&L and $I_t$ marking days beyond VaR. What is its mean if the model is right? If the true volatility is 1.2 times the model's? Where would you set the 5\% cut-off, and what power does it have?
\end{question}

\begin{question}[$\star\star\star$ \iqroles{Risk, QR} \iqfirm{any} \iqlevel{junior, experienced}]\label{iq:iv2:risk-validation-and-regulation:29}
A bank uses a 250-day historical simulation for a 99\% VaR. Volatility doubles overnight and stays doubled for 20 days. What is the exception rate on those days, and the expected count? How does an exponentially weighted model with $\lambda = 0.94$ do on the tenth day?
\end{question}

\begin{question}[$\star\star\star$ \iqroles{Risk, QR} \iqfirm{any} \iqlevel{experienced}]\label{iq:iv2:risk-validation-and-regulation:30}
Losses above a threshold of 2 follow a generalised Pareto distribution with shape 0.2 and scale 1, and 5\% of days exceed the threshold. What are the 99\% VaR and the 99\% expected shortfall? What happens as the shape reaches 1?
\end{question}

\begin{question}[$\star\star\star$ \iqroles{Risk, QR} \iqfirm{any} \iqlevel{experienced}]\label{iq:iv2:risk-validation-and-regulation:31}
Two assets have a correlation of 0.5. What is the chance that both fall below their own 1\% quantile, under a Gaussian copula and under a Student-$t$ copula with 4 degrees of freedom? What is the limit of the tail dependence in each?
\end{question}

\begin{question}[$\star\star\star$ \iqroles{Risk, QR} \iqfirm{bank} \iqlevel{experienced}]\label{iq:iv2:risk-validation-and-regulation:32}
Under the trading-book rules, a 97.5\% shortfall is computed at a 10-day base horizon and then cascaded by liquidity horizons of 10, 20, 40, 60 and 120 days. The 10-day shortfall on all factors is 10, and on the factors with horizon at least 20, 40, 60 and 120 days it is 8, 5, 3 and 2. What is the capital measure? Why is it not $10\sqrt{12}$?
\end{question}

\begin{question}[$\star\star\star$ \iqroles{Risk, QR, Dev} \iqfirm{any} \iqlevel{junior, experienced}]\label{iq:iv2:risk-validation-and-regulation:33}
A VaR spreadsheet computes a rate change as $(x_t - x_{t-1})/(x_t + x_{t-1})$ when it was meant to divide by the average of the two. What happens to the volatility and the VaR? How many exceptions a year would you expect in a backtest, for normal returns? How would a validator catch this before any backtest?
\end{question}

\begin{question}[$\star\star\star$ \iqroles{Risk, QR} \iqfirm{bank} \iqlevel{experienced}]\label{iq:iv2:risk-validation-and-regulation:34}
Exposure at default is $1.4\,(\mathrm{RC} + \mathrm{mult} \times \mathrm{AddOn})$, with $\mathrm{mult} = \min\bigl(1,\ 0.05 + 0.95\exp((V - C)/(2 \times 0.95 \times \mathrm{AddOn}))\bigr)$ and $\mathrm{RC} = \max(V - C, 0)$. Take an add-on of 10. What is it at $V - C = 5$? At $-20$? What is the floor?
\end{question}

\begin{question}[$\star\star\star$ \iqroles{Risk, QR} \iqfirm{bank, hedge fund} \iqlevel{experienced}]\label{iq:iv2:risk-validation-and-regulation:35}
A portfolio's P\&L is $2x_1 - x_2$ in millions, where the factors have mean 0, variances 4 and 1, and covariance 1. What is the most likely scenario that loses exactly 10? How many standard deviations away is it? Why does the second factor move the way it does?
\end{question}

\begin{question}[$\star\star\star$ \iqroles{Risk, QR} \iqfirm{any} \iqlevel{junior, experienced}]\label{iq:iv2:risk-validation-and-regulation:36}
A risk manager writes a stressed correlation matrix for three assets: correlations of 0.9 for pairs (1, 2) and (2, 3), and $-0.5$ for pair (1, 3). What is wrong with it? What range can the (1, 3) correlation take? What can go wrong if you use the matrix anyway, and how do you repair it?
\end{question}

\begin{question}[$\star\star\star$ \iqroles{Risk, QR} \iqfirm{bank} \iqlevel{experienced}]\label{iq:iv2:risk-validation-and-regulation:37}
You backtest a ten-day 99\% VaR on overlapping ten-day windows over 250 days. For independent daily normal returns, what are the exception count's mean and variance, against a binomial with independent days? What does this do to the usual test, and what do you do about it?
\end{question}

\begin{question}[$\star\star\star\star$ \iqroles{Risk, QR} \iqfirm{bank} \iqlevel{experienced}]\label{iq:iv2:risk-validation-and-regulation:38}
You are validating a desk's 99\% one-day VaR. In 500 days it had 9 exceptions. The transitions are: 485 no-no, 5 no-yes, 5 yes-no, 4 yes-yes. Run both tests. Then go through what you would look at next, and what you would tell the desk and the regulator.
\end{question}

\begin{question}[$\star\star\star\star$ \iqroles{Risk, Trader, QR} \iqfirm{hedge fund, bank} \iqlevel{experienced}]\label{iq:iv2:risk-validation-and-regulation:39}
An options book is delta-hedged, with gamma $-20{,}000$ shares per dollar, vega $-\$40{,}000$ per volatility point, and theta $+\$22{,}500$ a day. The stock's daily move is normal with a standard deviation of \$1.50. Implied volatility moves by a normal of 1 point a day, with correlation $-0.7$ to the stock. Estimate the 99\% VaR and shortfall: delta-normal, gamma alone, gamma and vega together. What happens with Student-$t$ moves of 4 degrees of freedom? What about a stress of $-\$10$ and $+5$ points? What limit do you set?
\end{question}

\begin{question}[$\star\star\star\star$ \iqroles{Risk, QR, QD} \iqfirm{bank} \iqlevel{experienced}]\label{iq:iv2:risk-validation-and-regulation:40}
A colleague's binomial tree for a European call gives these, with $S = K = 100$, $r = 5\%$, $\sigma = 20\%$ and $T = 1$. The closed form is 10.4506. At $N = 50$, 10.4107. At 51, 10.4850. At 100, 10.4306. At 101, 10.4680. At 200, 10.4406. At 201, 10.4593. At a strike of 180, $N = 50$ gives 0.0249 and $N = 200$ gives 0.0276, and the closed form is 0.0286. At $\sigma = 0.1\%$ and $N = 500$, it returns $2.7 \times 10^{159}$. What do you conclude? What is the fix for each? What else do you test before you sign?
\end{question}
